\section{Online Algorithm}
\label{sec:policy}

PACE's online compilation algorithm uses observed task arrivals and compilation outcomes to make cost-aware compilation proposals, with a separate cumulative budget to admit charges.
The proposal can be wrong about future arrivals or verification, while the budget uses only arrivals already observed and costs already charged.
The algorithm serves the current task before considering compilation for later tasks.

\subsection{Cost-aware compilation proposal}
\label{sec:fullpolicy}

A family becomes eligible after $k$ completed agent runs, including fallbacks.
Let $N$ be its observed arrival count and $d$ its age in total stream arrivals.
After $j$ compilation attempts and $g$ successful verifications, define
\begin{equation}
\hat H=\min\!\left(\frac{Nd}{d+5},\frac1h\right),\qquad
\hat p=\frac{g+1}{j+2},\qquad
\hat B=C+\left(\frac1{\hat p}-1\right)C^{\mathrm{fail}}.
\label{eq:estimated-price}
\end{equation}
We use $1/0=\infty$.
The smoothed count estimates future uses, capped by the supplied mean program lifetime.
It is a heuristic estimate from observed arrivals, without a fitted prediction model.
The estimated compilation cost $\hat B$ includes failed attempts before obtaining a verified program, under the independent-retry interpretation derived in Appendix~\ref{sec:breakeven}.
Neither interpretation is required for the cost guarantee.

With drift-adjusted saving $s_h=c-c^{\mathrm{prog}}-[h+(1-h)q]c$, an eligible family without a live program proposes compilation when
\begin{equation}
s_h>0\quad\text{and}\quad\hat Hs_h>\hat B+m\min(T,d).
\label{eq:trigger}
\end{equation}
The final term allows for one manifest entry during its estimated residence.
Entries expire after $T$ idle arrivals, while stored programs can be relisted until a detected break invalidates them.
The proposal accounts for the two sources of uncertain payback separately.
Low verification probability increases the estimated number of paid failures before obtaining a program, whereas GUI drift reduces the estimated number of uses after verification.
An expensive failed attempt can therefore make compilation unattractive even when successful construction is cheap.
The following budget checks control the cost incurred when these estimates are wrong.
Appendix~\ref{sec:matched-budget} compares alternative proposals with the same budget.
Their similar costs do not identify the exact smoothing formula as a separate contribution.

\subsection{Budget checks}
\label{sec:cost-protection}

Before routing, PACE expires idle entries and tentatively relists any stored live program for the arriving family.
The manifest exposes programs to the model for routing, and its cost is paid on every arrival, including tasks served by the agent.
Storing an artifact incurs no token cost, but leaving it listed can keep incurring routing cost.
For the resulting manifest size $\ell_t$, it reserves
\begin{equation}
U_t^{\mathrm{serv}}=\tau_0+m\ell_t+
\begin{cases}
c,&\text{agent service},\\
c^{\mathrm{prog}}+c,&\text{program service}.
\end{cases}
\label{eq:service-reserve}
\end{equation}
The program reservation includes a possible fallback.
Service proceeds only if $K_{t-1}+U_t^{\mathrm{serv}}\leq(1+\epsilon)A_t$.
Otherwise PACE clears the manifest before routing and uses the agent, whose charge is at most $a_t$.
Clearing preserves stored artifacts and leaves no mandatory future charge.

Let $\widetilde K_t$ be the cumulative charge after service.
A proposed compilation proceeds only if
\begin{equation}
\widetilde K_t+\max(C,C^{\mathrm{fail}})\leq(1+\epsilon)A_t.
\label{eq:compile-admission}
\end{equation}
PACE then adds the realized compilation charge to its cumulative cost, including failed attempts and any compilation after the last service.
Savings can fund compilation across families, but estimated future uses add no budget.
Appendix~\ref{app:protocol-details} gives the complete execution order.

\paragraph{Budget example.}
Consider a proposed compilation with success and failure charges of $60$ and $120$ weighted tokens.
If the cumulative reference and cost already charged are both $300$, a margin of $0.25$ leaves $75$ tokens available.
PACE defers compilation because a failure would exceed the $375$ limit, even if the estimated saving justifies its estimated compilation cost.
After two further agent-served arrivals costing $100$ each, the reference and charged cost are both $500$.
The available $125$ tokens now cover either outcome.
Only observed arrivals increase this allowance.

\subsection{Cost guarantee}
\label{sec:guarantee}

\begin{theorem}[Cost bound relative to agent execution]
\label{thm:cost-protection}
Let the reference increments $a_t$ be nonnegative.
Suppose every proposed action has a known finite charge bound, and clearing the manifest before routing costs no tokens, leaves no mandatory future charge, and permits default service at cost at most $a_t$.
For any $\epsilon\geq0$, PACE with exact arithmetic and valid reservations satisfies
\begin{equation}
K_t\leq(1+\epsilon)A_t
\quad\text{for every completed-arrival prefix }t
\label{eq:cost-guarantee}
\end{equation}
on every history allowed by these charge bounds.
\end{theorem}

\begin{proof}
The bound is preserved in two steps.
An admitted service costs no more than its reservation.
If admission fails, default service adds at most $a_t$, so
\[
K_{t-1}+a_t\leq(1+\epsilon)A_{t-1}+a_t\leq(1+\epsilon)A_t.
\]
The compilation check then preserves the same bound for either outcome.
Starting from $K_0=A_0=0$ establishes the result by induction.
\end{proof}

The fixed-charge model satisfies the assumptions with $a_t=\tau_0+c$.
The implementation permits tolerance $10^{-8}\max\{1,(1+\epsilon)A_t\}$.
The bound is global over the stream and concerns cost, not task errors or per-family spending.
For variable real bills, reservations must be enforceable and default service affordable.
Historical means alone do not supply these conditions.
PACE is a feasible algorithm for Equation~\ref{eq:objective}, without a claim of optimizing it.
